% How to read Julia listings in this manuscript.
% Included from the appendix of main.tex.

\section{How to read the Julia listings}
\label{app:code-listings}

Code examples throughout this paper are displayed in framed Julia listings.
A blue \textbf{Julia} badge in the top-right corner identifies the language.
Line numbers on the left are the ones referred to in the text
(for example, Listing~\ref{lst:listing-legend}, lines~4--6).
Syntax highlighting is chosen to make the origin of each identifier visible at a
glance, rather than to reproduce a particular editor theme.

The colour roles are as follows.
\begin{itemize}
    \item \textbf{Blue, bold:} Julia keywords such as \jlinline{using}.
    \item \textbf{Violet, bold:} names exported by \texttt{ProcessTensors.jl}.
    This includes package-specific constructors and converters
    (\jlinline{to_dm}, \jlinline{liouv_sites}, \jlinline{liouvillian_mpo},
    \ldots) as well as the tensor-network types that form the package's public
    surface, notably \jlinline{MPS}, \jlinline{MPO}, and \jlinline{OpSum}.
    \item \textbf{Teal, bold:} functions that belong to \texttt{ITensors.jl}
    and are \emph{not} part of the \texttt{ProcessTensors.jl} export list
    (\jlinline{Index}, \jlinline{ITensor}, \jlinline{siteinds},
    \jlinline{dag}, \jlinline{prime}, \jlinline{delta}, \jlinline{inner},
    \jlinline{op}, \jlinline{apply}, \ldots).
    \item \textbf{Orange, bold:} package types and backends such as
    \jlinline{Hilbert}, \jlinline{Liouville}, \jlinline{Dense}, and
    \jlinline{ACE}.
    \item \textbf{Near-black:} ordinary variables.
    \item \textbf{Green, bold:} numeric literals.
    \item \textbf{Crimson:} string literals.
    \item \textbf{Grey italic:} comments.
\end{itemize}

Listing~\ref{lst:listing-legend} collects these conventions in one place.
Violet names are the objects a user of \texttt{ProcessTensors.jl} is expected
to call; teal names are the underlying ITensor primitives on which those
objects are built.

\begin{juliacode}[
  caption={Colour legend for the listings in this paper.},
  label={lst:listing-legend},
]
using ITensors, ProcessTensors

s = Index(2, "s")                       # ITensor index
psi = ITensor([1.0, 0.0], s)
sites = siteinds("S=1/2", 1)

rho = to_dm(MPS(sites, ["Up"]))         # ProcessTensors converter + MPS
H = OpSum()
H += 0.7, "Sx", 1
L = liouvillian_mpo(H, liouv_sites(sites))
\end{juliacode}
